\section{似然比比较两个模型对数据的解释力}
\label{sec:11-Mathematical-Statistics/05-Hypothesis-Testing/03}

线上要加一个漂移报警：每小时算一批特征的统计量，判断当前流量是照常还是已经跑偏。历史数据两样都有——正常时段的分布 \(p_0\)，几次事故时段的分布 \(p_1\)。运维给的约束很硬：误报每周不超过一次，在这个前提下越早发现越好。

报警规则该按什么量排序？候选很多：偏离均值多少个标准差、几个特征同时越界、最大分量超没超过某个数。上一节结尾留下的问题就是它——同一份数据能造出无穷多个水平 \(\alpha\) 的规则，它们的功效相差很远。这一节给出选择的依据，依据的形状是一个比值。

\subsection{两个模型在同一个观测点上比高度}

先把问题削到最干净的形状：两个假设都完全指定了分布，

\[
H_0:X\sim p_0,\qquad H_1:X\sim p_1 .
\]

观测到 \(x\) 之后，两个模型各自给它分配了一个密度值 \(p_0(x)\) 与 \(p_1(x)\)。数据对两者的相对支持由似然比衡量：

\[
\Lambda_{10}(x)=\frac{p_1(x)}{p_0(x)} .
\]

抛硬币十次得到八次正面，公平硬币给这份数据的概率是 \(\binom{10}{8}0.5^{10}\approx 0.044\)，\(p=0.7\) 的硬币给的是 \(\binom{10}{8}0.7^{8}0.3^{2}\approx 0.233\)，比值 \(5.3\)。数据在第二个模型手里更容易出现。

要和上一节划清界限。\(p\) 值算的是 \(H_0\) 下的尾部概率，问的是``\(H_0\) 成立时看到这么极端的数据有多意外''；似然比是两个密度在同一个观测点上的高度之比，问的是``这一个点，谁给它的份额更大''。一个沿着尾部积分，一个在单点取值，两个量指向不同的问题，数值上也没有换算关系。

比值大于 \(1\) 就报警显然不行：\(H_0\) 为真时，纯靠涨落，比值也有相当的概率越过 \(1\)。需要的是一个阈值 \(c\)，它把误报率钉在 \(\alpha\)，同时让漏报尽可能少。Neyman 与 Pearson 在 1930 年代证明了这个阈值规则的最优性，结论是：不存在比它更好的检验。

\subsection{固定预算下的装箱问题}

把样本空间切成很多小块。第 \(j\) 块进入拒绝域，成本是它在 \(H_0\) 下的概率质量 \(w_j=p_0(x_j)\Delta x\)，收益是它在 \(H_1\) 下的概率质量 \(p_1(x_j)\Delta x=w_j\cdot\Lambda_{10}(x_j)\)。所有被选中块的成本之和不得超过 \(\alpha\)，收益之和就是功效。

\begin{center}
\begin{tikzpicture}[x=0.78cm,y=0.85cm,font=\footnotesize,>=stealth]
  \foreach \a/\w/\h in {0/0.45/3.10, 0.50/0.55/2.55, 1.10/0.60/2.05}
    \filldraw[fill=black!35,draw=black!60] (\a,0) rectangle (\a+\w,\h);
  \foreach \a/\w/\h in {1.75/0.75/1.60, 2.55/0.85/1.20, 3.45/0.95/0.88,
                        4.45/1.05/0.62, 5.55/1.15/0.40, 6.75/1.25/0.22, 8.05/1.35/0.10}
    \filldraw[fill=black!8,draw=black!45] (\a,0) rectangle (\a+\w,\h);
  \draw[black!70] (-0.15,0) -- (9.6,0);
  \draw[black!70] (-0.15,0) -- (-0.15,3.35);
  \draw[densely dashed,black!85] (-0.15,1.83) -- (9.6,1.83);
  \node[anchor=west,black!85] at (9.7,1.83) {\(c\)};
  \node[anchor=east,black!70,rotate=90] at (-0.55,2.4) {似然比 \(p_1/p_0\)};
  \node[anchor=north,black!70] at (4.7,-0.35) {样本空间按似然比从高到低排列，宽度 \(=\) 在 \(H_0\) 下的概率质量};
  \draw[<->,black!70] (0,3.45) -- (1.70,3.45);
  \node[anchor=west,black!70] at (1.85,3.45) {阴影宽度合计 \(=\alpha\)，阴影面积 \(=\) 功效};
\end{tikzpicture}

\emph{一条水平线 \(c\) 就是全部的自由度：线上方的块进入拒绝域，宽度受 \(\alpha\) 约束，面积（宽 \(\times\) 高）等于功效。}
\end{center}

图里两个量的对应值得停一下：宽度是 \(H_0\) 下的质量，高度是似然比，宽乘高恰好是 \(H_1\) 下的质量。于是``在总宽度不超过 \(\alpha\) 的前提下让总面积最大''这句话，字面意思就是``在第一类错误不超过 \(\alpha\) 的前提下让功效最大''。装箱问题的解一眼可见：从最高的块开始拿，一直拿到宽度用完。而``最高的块''按定义就是似然比最大的那些点，它们围成的区域正是 \(\set{x:p_1(x)/p_0(x)>c}\)。

Neyman--Pearson 引理把这句话写成定理：在所有 size 不超过 \(\alpha\) 的检验中，拒绝域取

\[
\frac{p_1(X)}{p_0(X)}>c ,
\]

并在等号处适当随机化（离散分布下为了把 size 精确凑到 \(\alpha\)），对指定的简单备择 \(H_1\) 获得最大功效。

反过来看更能说明为什么别无选择。若某个替代检验的 size 相同而功效更高，那么它的拒绝域必定包含了一些比值较低的块，同时漏掉了一些比值较高的块；把这两部分对调，宽度不变而面积增加，与它的最优性矛盾。严格证明把这段交换论证写成一个非负积分不等式，占不到一页，此处不展开。

结论的适用范围要说清楚。最优性是针对一个具体的简单备择而言的。备择是一整族分布时，对 \(\theta_1\) 最优的检验对 \(\theta_2\) 可能很差，因为 \(p_{\theta_1}/p_0\) 与 \(p_{\theta_2}/p_0\) 给样本空间的排序未必一致。一致最优检验只在似然比对某个统计量单调的结构里存在——指数族是这样的结构。

\subsection{单调性把拒绝域收缩到一侧}

看正态的情形。\(X\sim N(\mu,\sigma^2)\)，检验 \(\mu=\mu_0\) 对 \(\mu=\mu_1>\mu_0\)，对数似然比是

\[
\log\frac{p_{\mu_1}(X)}{p_{\mu_0}(X)}=\frac{\mu_1-\mu_0}{\sigma^{2}}X+\text{与 }X\text{ 无关的项} .
\]

它随 \(X\) 单调递增，所以``似然比超过 \(c\)''与``\(X\) 超过某个数''是同一个集合。上一节那个单侧 \(z\) 检验的形状，到这里才有了来历：拒绝域取右尾，并非出于``大值可疑''的直觉，而是似然比的排序把它逼到了那里。同时可以看出，只要 \(\mu_1>\mu_0\)，具体取多大都不影响拒绝域的形状，仅仅改变 \(c\) 与 \(\alpha\) 的对应——所以对整个单侧备择族，这个检验一致最优。

一般的指数族密度写成 \(\exp\set{\eta(\theta)T(x)-A(\theta)}h(x)\)，两个参数值的对数似然比是 \(T(x)\) 的线性函数，\(\eta\) 单调时立刻得到单调似然比。\(t\) 检验、\(F\) 检验、Poisson 与 Binomial 检验的拒绝域方向，来源都在这里。

双侧备择就没这么幸运。\(\mu\ne\mu_0\) 里同时有大于和小于两簇，对一簇最优的检验拒绝大值，对另一簇最优的拒绝小值，不存在对两边同时最优的单一拒绝域。开头那个漂移报警也是这个处境：漂移可能往任何方向去，\(p_1\) 不是一个分布而是一族。需要一个不挑方向的构造。

\subsection{让每个模型先各自拟合到最好，再比峰高}

若 \(H_0\) 把参数限制在 \(\Theta_0\) 内，完整模型允许参数在 \(\Theta\) 内自由取值，广义似然比定义为

\[
\Lambda(X)=\frac{\sup_{\theta\in\Theta_0}L(\theta;X)}{\sup_{\theta\in\Theta}L(\theta;X)} .
\]

分子是绑着手的最佳拟合，分母是放开手的最佳拟合，\(0\le\Lambda\le 1\)。约束下的最优拟合与自由拟合差得远，\(\Lambda\) 就小——数据在说，即使在你划定的范围内挑遍所有参数，也解释不了我。两个 \(\sup\) 都由最大似然给出，用法见\hyperref[sec:11-Mathematical-Statistics/02-Point-Estimation/03]{``最大似然选择最能解释数据的参数''}。

直接用 \(\Lambda\) 不方便，因为没有现成的分布表。习惯上取

\[
-2\log\Lambda ,
\]

一来把取值从 \((0,1]\) 搬到 \([0,\infty)\)，二来 Wilks 证明了在正则条件下

\[
-2\log\Lambda\ \overset{d}{\longrightarrow}\ \chi^{2}_{r},
\]

其中 \(r\) 是完整模型与受限模型的自由参数个数之差。骨架两步：在真参数附近把对数似然展成二次型，最大值之差化为 score 的二次型；再用中心极限定理，score 渐近正态，其二次型服从卡方，自由度等于被约束掉的方向数。条件对账：真参数落在参数空间内部、模型可识别、Fisher 信息非奇异、参数维度不随样本量增长、似然二阶可微。这几条里任何一条失效，自由度规则都会跟着失效。

嵌套模型比较是它最常见的用法。两组数据各拟合一条回归线是四个参数，假设两组共用一条线是两个参数，\(r=2\)，算出 \(-2\log\Lambda=6.3\) 查 \(\chi^{2}_{2}\) 得 \(p\approx 0.043\)。特征选择里``加这三个特征值不值''、混合模型里``多一个成分有没有必要''，都是同一个模板。不必为每种比较重新推导分布。

\subsection{三种检验站在似然曲面的三个位置}

同一片似然曲面可以从三处探测。Wald 检验站在无约束的极大似然估计处，用估计的曲率去量它离约束面有多远，只需拟合完整模型一次；麻烦是结果依赖参数化方式，把 \(\theta\) 换成 \(\log\theta\) 统计量就变了，由此得到的置信区间偶尔会越出参数的定义域。Score 检验站在受限拟合处，检查那一点的似然梯度是否接近零，梯度仍陡说明数据在往外拉；它只需拟合受限模型，计算量最小，拟合优度里的 Pearson 卡方检验就是它在多项分布下的特例。似然比检验比较两个峰顶的高度差，要拟合两个模型，计算量最大，但对参数变换不变，有限样本下通常最接近标称水平。

三者在正则条件下渐近等价，差别在似然曲面偏离二次型的程度。曲面越对称，三次测量越一致；参数靠近边界、样本偏小、曲率变化剧烈时，似然比检验最稳，Wald 最先出问题。

\subsection{卡方近似失效的两个典型场合}

真参数落在参数空间边界上是头一个。检验随机效应的方差是否为零，而方差不能为负，真值 \(0\) 恰在边界；最大似然估计有正概率堆在边界上，二次展开不成立，极限分布变成卡方的混合，例如 \(0.5\chi^{2}_{0}+0.5\chi^{2}_{1}\)。照常查表会得到偏保守的 \(p\) 值。

参数在原假设下不可识别是第二个。检验混合模型的成分个数时，``一个成分就够''这个原假设让多余成分的位置参数完全不影响似然，信息矩阵奇异，自由度无从谈起。参数维度随样本量增长、模型设定错误、工具变量很弱，也都会让渐近结论落空。

补救办法通常是参数 bootstrap：从拟合好的受限模型里重复模拟数据，每次重算 \(-2\log\Lambda\)，用经验分布代替卡方表。计算贵，但它直接给出统计量在你自己这个模型下的分布，不再依赖任何渐近条件。

\subsection{比值本身不是概率，还差一个先验}

最后堵一个口子。似然比只衡量数据带来的更新力量，Bayes 的后验赔率还要乘上先验赔率：

\[
\frac{P(H_1\given x)}{P(H_0\given x)}
=\frac{p(x\given H_1)}{p(x\given H_0)}\cdot\frac{P(H_1)}{P(H_0)} .
\]

似然比是右端第一个因子。若 \(H_1\) 的先验赔率极低——筛查一种发病率百万分之一的疾病，或者从十万个特征里找那少数几个真有信号的——似然比等于 \(100\) 也只是把万分之一变成百分之一。这与上一节那张基准率的图说的是同一件事：数据这一侧的量再漂亮，也补不上先验那一侧的缺口。

复合模型下的 Bayes factor 对参数先验积分，而广义似然比取上确界。前者问``把参数空间里每个值按先验加权，整个模型得到多少支持''，多余的自由度会因为先验把权重摊到低似然区域而自动受罚；后者问``在你允许的范围里，拟合得最好的那个参数有多好''，不含这项惩罚，所以复合假设下它总是偏向参数更多的那个模型，这也是为什么它要配一张卡方表来兑现自由度的账。

到此为止，一次比较的最优形状已经清楚。麻烦在于现实里几乎没有人只做一次比较：一份实验报告里躺着十几个指标，一次特征筛选要过上万列。检验做很多次时，前面所有关于 \(\alpha\) 的承诺都要重新计价，这是下一节的事。
